\chapter{基于数据本身的知识获取}
\label{chap:data-derived-knowledge-acquisition}

前几章讨论的训练依据具有明确的任务含义：类别标签回答样本属于哪一类，数值标注回答目标量是多少，成对判断则回答哪一个候选更符合要求。可是，大量文本、图像和传感器记录只有观测内容，并没有随记录附带最终任务的答案。若等到每条记录都获得任务标签后才开始学习，未标注数据中已经存在的结构就没有得到利用。于是，本章要回答一个更基础的问题：没有现成任务标签时，怎样从数据中建立可训练、可检查的依据？

关键不在于目标函数是否写成“预测”，而在于预测目标从哪里来。遮住图像块后恢复像素，答案在遮蔽之前已经记录在数据中；根据一个裁剪预测另一个裁剪的表示，目标向量由训练中的模型计算；聚类后预测簇编号，目标又来自分组算法。三者都能写成输入到目标的学习问题，却依赖不同假设，也需要不同的技术检查。把它们统称为预测任务，会掩盖监督来源、梯度路径和失效原因。

本章按训练依据的来源组织内容。基于\term{观测内容}的方法从完整记录中保留待预测部分；基于\term{样本关系}的方法说明哪些观测应共享信息、哪些观测可作相对区分；基于\term{数据分组}的方法先发现群组，再把组编号或原型分配变成目标。上下文、时间顺序和多模态配对并不是第四种来源，它们分别规定观测或关系怎样组织。读者需要始终沿“已有信息、构造操作、训练目标、参数更新、下游使用”五个位置追踪一个方法。

图~\ref{fig:data-derived-knowledge-target-sources}先对照最容易混淆的两类目标。左侧的像素目标在模型计算之前就能从原记录中取出；右侧的表示目标只有在视图经过编码分支之后才出现，其数值会随参数变化。表示目标究竟由共享分支、停止梯度分支还是移动平均分支产生，要到第~\ref{sec:data-derived-knowledge-training-objectives}节结合具体方法判断。

\begin{figure*}[htbp]
  \centering
  \begin{tikzpicture}[
  x=1mm,y=1mm,
  font=\figurefont\footnotesize,
  panel/.style={draw=BookRule,fill=NNPanel,line width=.85pt,rounded corners=1.5mm},
  data/.style={draw=FigureInput,fill=FigureInput!12,line width=1.1pt,minimum height=10mm,align=center,inner sep=2mm},
  model/.style={draw=FigureModel,fill=FigureModel!12,line width=1.1pt,minimum height=10mm,align=center,inner sep=2mm},
  target/.style={draw=FigureOutput,fill=FigureOutput!10,line width=1.1pt,minimum height=10mm,align=center,inner sep=2mm},
  loss/.style={draw=FigureLoss,fill=FigureLoss!12,line width=1.1pt,minimum height=10mm,align=center,inner sep=2mm},
  flow/.style={knowledge arrow,draw=NNWire,line width=1.2pt,rounded corners=1mm},
  keep/.style={knowledge arrow,draw=FigureOutput,line width=1.1pt,rounded corners=1mm},
  annotation/.style={font=\figurefont\scriptsize,text=BookMuted,align=center}
]
  \draw[panel] (1,-76) rectangle (81,12);
  \draw[panel] (87,-76) rectangle (168,12);
  \node[font=\figurefont\bfseries\footnotesize,text=BookInk] at (41,7) {(a) 观测目标：答案随记录保留};
  \node[data,minimum width=24mm] (obs) at (41,-9) {完整观测\\$\bm{x}$};
  \node[data,minimum width=24mm] (visible) at (20,-28) {可见输入\\$\bm{x}_{V}$};
  \node[target,minimum width=24mm] (held) at (62,-28) {保留目标\\$\bm{x}_{M}$};
  \node[model,minimum width=24mm] (predict) at (20,-48) {模型预测\\$\widehat{\bm{x}}_{M}$};
  \node[loss,minimum width=24mm] (recloss) at (62,-48) {重建损失};
  \coordinate (obs-split) at (41,-18);
  \draw[NNWire,line width=1.2pt] (obs.south) -- (obs-split);
  \draw[flow] (obs-split) -| (visible.north);
  \draw[keep] (obs-split) -| (held.north);
  \draw[flow] (visible.south) -- (predict.north);
  \draw[flow] (predict.east) -- (recloss.west);
  \draw[keep] (held.south) -- (recloss.north);
  \node[annotation,text width=68mm] at (41,-66) {目标从原记录中预先保留，\\不需要编码分支生成答案。};

  \node[font=\figurefont\bfseries\footnotesize,text=BookInk] at (127.5,7) {(b) 表示目标：答案在训练中计算};
  \node[data,minimum width=24mm] (raw) at (127.5,-9) {原样本\\$\bm{x}$};
  \node[data,minimum width=24mm] (viewa) at (106,-28) {视图A\\$\bm{v}$};
  \node[data,minimum width=24mm] (viewb) at (149,-28) {视图B\\$\bm{v}'$};
  \node[model,minimum width=24mm] (enca) at (106,-48) {编码分支\\$\bm{z}$};
  \node[model,minimum width=24mm] (encb) at (149,-48) {编码分支\\$\bm{z}'$};
  \node[loss,minimum width=22mm] (match) at (127.5,-66) {匹配目标};
  \coordinate (view-split) at (127.5,-18);
  \draw[NNWire,line width=1.2pt] (raw.south) -- (view-split);
  \draw[flow] (view-split) -| (viewa.north);
  \draw[flow] (view-split) -| (viewb.north);
  \draw[flow] (viewa.south) -- (enca.north);
  \draw[flow] (viewb.south) -- (encb.north);
  \draw[flow] (enca.south) |- (match.west);
  \draw[flow] (encb.south) |- (match.east);

  \node[annotation,anchor=north west,text width=166mm,align=left] at (1,-80) {右侧表示在编码之后产生，随参数和训练步变化；本图只比较目标来源，不预设教师或停止梯度。};
\end{tikzpicture}

  \caption{观测目标与学得表示目标的来源。左侧从完整观测中预先保留答案，再让模型只读取可见部分；右侧先建立两视图的对应关系，再由编码分支在训练中计算要匹配的表示。图中没有预设右侧必须使用教师或停止梯度。}
  \label{fig:data-derived-knowledge-target-sources}
\end{figure*}
\FloatBarrier

\section{基于观测内容的知识获取}
\label{sec:data-derived-knowledge-observation}

一条完整记录常常可以拆成当前允许模型读取的部分和暂时不允许读取的部分。后者虽然在构造训练样本时已经存在，却可以被保留下来作为答案。\term{观测目标}就是这样从已记录的词元、像素、属性或未来片段中取出的预测对象。它不需要正在学习的编码器来生成，但仍依赖遮蔽位置、预测方向、预处理方式和损失范围等人为选择。

设完整观测为$\bm{x}=(x_j)_{j\in I}$，$I$是位置集合。构造过程选择可见集合$V\subset I$和目标集合$M\subseteq I$，通常要求模型输入中不能直接恢复出$\bm{x}_M$的原值。模型读取处理后的$\widetilde{\bm{x}}_V$，输出$\widehat{\bm{x}}_M$或目标的条件分布。这里要分别记录完整观测$\bm{x}$、可见输入$\widetilde{\bm{x}}_V$、保留目标$\bm{x}_M$和损失作用位置$M$；只保存一份“训练样本”而不保留这些角色，难以排查信息泄露。

表~\ref{tab:data-derived-knowledge-observation-tasks}比较四种常见安排。它们可以组合：时序记录中的一段未来也可以先被扰动，再要求模型恢复。划分依据不是损失名称，而是观测怎样被组织成输入与答案。

\begin{table*}[htbp]
  \centering
  \renewcommand{\arraystretch}{1.15}
  \begin{tabularx}{\linewidth}{@{}>{\RaggedRight\setlength{\parindent}{0pt}\arraybackslash}p{22mm}*{3}{>{\RaggedRight\setlength{\parindent}{0pt}\arraybackslash}X}@{}}
    \toprule
    构造方式 & 模型可见部分 & 保留目标 & 损失位置 \\
    \midrule
    掩码恢复 & 隐藏选中图像块后的图像 & 被遮蔽块的原像素 & 通常只在掩码位置 \\
    去噪恢复 & 加入扰动后的图像 & 扰动前的原始记录 & 受扰位置或全部位置 \\
    上下文预测 & 前缀“花朵逐渐” & 后续词元“开放” & 约定的预测位置 \\
    时序预测 & 判断时点以前的历史窗口 & 后续帧或状态 & 一个未来点或一段轨迹 \\
    \bottomrule
  \end{tabularx}
  \caption{观测预测任务中的可见信息与目标。未来片段在构造训练集时已经记录，但在真正作出预测的时点尚不可见；四行比较输入目标安排，不表示这些构造互斥。}
  \label{tab:data-derived-knowledge-observation-tasks}
\end{table*}
\FloatBarrier

\subsection{基于遮蔽恢复的目标构造}
\label{sec:data-derived-knowledge-masking-denoising}

遮蔽恢复与去噪恢复都从同一完整样本产生输入目标对。前者直接隐藏若干单位，后者保留输入位置但修改其内容；遮蔽也可以视为把选中位置替换成特殊缺失符号的一种扰动。将两者分开讨论，是为了看清可见信息和损失位置，而不是声称它们在理论上互斥。

图~\ref{fig:data-derived-knowledge-masking-denoising}使用同一$3\times3$数值网格说明差别。上行让两个位置完全不可见，下行则把同样位置的数值改动。赭色边框只标出这个教学例子参与损失计算的位置。真实任务可以在全部位置计算去噪损失，也可以只在已知被扰动的位置计算；这一选择必须写进目标定义。

\begin{figure*}[htbp]
  \centering
  \begin{tikzpicture}[
  x=1mm,y=1mm,
  font=\figurefont\scriptsize,
  flow/.style={knowledge arrow,draw=NNWire,line width=1.1pt},
  grid/.style={draw=BookBlue,line width=.85pt},
  cell/.style={draw=BookRule,line width=.5pt,minimum size=5mm,inner sep=0pt},
  heading/.style={font=\figurefont\bfseries\footnotesize,text=BookInk},
  stage/.style={font=\figurefont\footnotesize,text=BookMuted}
]
  \node[heading,anchor=east] at (13,1) {掩码};
  \node[heading,anchor=east] at (13,-34) {去噪};
  \foreach \x/\name in {23/原始观测,64/处理后输入,105/模型预测,146/保留目标}{
    \node[stage] at (\x,13) {\name};
  }

  \begin{scope}[shift={(15,8)}]
    \foreach \r/\c/\v in {0/0/.2,0/1/.4,0/2/.6,1/0/.3,1/1/.8,1/2/.7,2/0/.1,2/1/.5,2/2/.9}{
      \node[cell] at ({\c*5},{-\r*5}) {$\v$};
    }
  \end{scope}
  \begin{scope}[shift={(56,8)}]
    \foreach \r/\c/\v in {0/0/.2,0/1/.4,0/2/{\times},1/0/.3,1/1/{\times},1/2/.7,2/0/.1,2/1/.5,2/2/.9}{
      \node[cell] at ({\c*5},{-\r*5}) {$\v$};
    }
    \fill[FigureOutput!18] (7.5,-2.5) rectangle (12.5,2.5);
    \fill[FigureOutput!18] (2.5,-7.5) rectangle (7.5,-2.5);
    \node[cell] at (10,0) {$\times$};
    \node[cell] at (5,-5) {$\times$};
  \end{scope}
  \begin{scope}[shift={(97,8)}]
    \foreach \r/\c/\v in {0/0/.2,0/1/.4,0/2/.55,1/0/.3,1/1/.70,1/2/.7,2/0/.1,2/1/.5,2/2/.9}{
      \node[cell] at ({\c*5},{-\r*5}) {$\v$};
    }
    \draw[FigureLoss,line width=1.1pt] (7.5,-2.5) rectangle (12.5,2.5);
    \draw[FigureLoss,line width=1.1pt] (2.5,-7.5) rectangle (7.5,-2.5);
  \end{scope}
  \begin{scope}[shift={(138,8)}]
    \foreach \r/\c/\v in {0/0/.2,0/1/.4,0/2/.6,1/0/.3,1/1/.8,1/2/.7,2/0/.1,2/1/.5,2/2/.9}{
      \node[cell] at ({\c*5},{-\r*5}) {$\v$};
    }
    \draw[FigureLoss,line width=1.1pt] (7.5,-2.5) rectangle (12.5,2.5);
    \draw[FigureLoss,line width=1.1pt] (2.5,-7.5) rectangle (7.5,-2.5);
  \end{scope}
  \draw[flow] (31,3) -- (53,3);
  \draw[flow] (72,3) -- (94,3);
  \draw[flow] (113,3) -- node[above,text=FigureLoss] {仅比较框内位置} (135,3);

  \begin{scope}[shift={(15,-27)}]
    \foreach \r/\c/\v in {0/0/.2,0/1/.4,0/2/.6,1/0/.3,1/1/.8,1/2/.7,2/0/.1,2/1/.5,2/2/.9}{
      \node[cell] at ({\c*5},{-\r*5}) {$\v$};
    }
  \end{scope}
  \begin{scope}[shift={(56,-27)}]
    \foreach \r/\c/\v in {0/0/.2,0/1/.4,0/2/.8,1/0/.3,1/1/.6,1/2/.7,2/0/.1,2/1/.5,2/2/.9}{
      \node[cell] at ({\c*5},{-\r*5}) {$\v$};
    }
    \fill[BookAmber!18] (7.5,-2.5) rectangle (12.5,2.5);
    \fill[BookAmber!18] (2.5,-7.5) rectangle (7.5,-2.5);
    \node[cell] at (10,0) {.8};
    \node[cell] at (5,-5) {.6};
  \end{scope}
  \begin{scope}[shift={(97,-27)}]
    \foreach \r/\c/\v in {0/0/.2,0/1/.4,0/2/.65,1/0/.3,1/1/.72,1/2/.7,2/0/.1,2/1/.5,2/2/.9}{
      \node[cell] at ({\c*5},{-\r*5}) {$\v$};
    }
    \draw[FigureLoss,line width=1.1pt] (7.5,-2.5) rectangle (12.5,2.5);
    \draw[FigureLoss,line width=1.1pt] (2.5,-7.5) rectangle (7.5,-2.5);
  \end{scope}
  \begin{scope}[shift={(138,-27)}]
    \foreach \r/\c/\v in {0/0/.2,0/1/.4,0/2/.6,1/0/.3,1/1/.8,1/2/.7,2/0/.1,2/1/.5,2/2/.9}{
      \node[cell] at ({\c*5},{-\r*5}) {$\v$};
    }
    \draw[FigureLoss,line width=1.1pt] (7.5,-2.5) rectangle (12.5,2.5);
    \draw[FigureLoss,line width=1.1pt] (2.5,-7.5) rectangle (7.5,-2.5);
  \end{scope}
  \draw[flow] (31,-32) -- (53,-32);
  \draw[flow] (72,-32) -- (94,-32);
  \draw[flow] (113,-32) -- node[above,text=FigureLoss] {此例比较受扰位置} (135,-32);

  \node[font=\figurefont\scriptsize,text=BookMuted,anchor=west] at (15,-48) {数值与扰动均为教学构造；保留目标是原始记录，不代表无采集噪声的真值。};
\end{tikzpicture}

  \caption{遮蔽与去噪怎样构造输入目标对。数值网格和扰动是便于复核的教学构造：掩码位置用叉号表示，受扰位置用浅赭色表示，赭色边框表示本例计算损失的位置。保留目标来自扰动前的原始记录，但原始记录本身仍可能含有采集噪声。}
  \label{fig:data-derived-knowledge-masking-denoising}
\end{figure*}
\FloatBarrier

\subsubsection{基于掩码的内容恢复}
\label{sec:data-derived-knowledge-masked-recovery}

\term{掩码恢复}（masked reconstruction）先选择待恢复位置$M$，再让模型只根据其余位置$V=I\setminus M$预测$\bm{x}_M$。文本可以按词元隐藏，图像可以按不重叠图像块隐藏，表格则可以按字段或单元格隐藏。隐藏单位决定了模型需要跨越多大范围寻找信息：遮住单个像素常可依靠邻域插值，遮住连续区域或完整属性则需要更远的上下文。

若目标是连续像素，常见的掩码均方误差为
\begin{equation}
  \mathcal{L}_{\mathrm{mask}}
  =\frac{1}{|M|}\sum_{j\in M}\lVert \widehat{\bm{x}}_j-\bm{x}_j\rVert_2^2.
  \label{eq:data-derived-knowledge-masked-reconstruction}
\end{equation}
若目标是离散词元，则把每个位置的平方误差换成交叉熵。分母按被遮蔽单位数归一化，使不同$|M|$的样本具有可比尺度。损失若还覆盖未遮蔽位置，模型可以通过复制大量可见内容降低平均值；这不一定错误，却改变了目标中“恢复未知内容”和“复现已知内容”的相对权重。

继续用图~\ref{fig:data-derived-knowledge-masking-denoising}上行的两个位置。假设保留目标为$(0.6,0.8)$，模型预测为$(0.55,0.70)$，则式~\eqref{eq:data-derived-knowledge-masked-reconstruction}给出
\[
  \mathcal{L}_{\mathrm{mask}}
  =\frac{(0.55-0.6)^2+(0.70-0.8)^2}{2}
  =0.00625.
\]
这个数值只检查两处预测，不把其余七个已经可见的值计入成绩。例子也没有声称$0.6$和$0.8$是真实世界的无噪声像素，它们只是训练记录中被暂时保留的值。

掩码自编码器（masked autoencoder, MAE）把图像划分为块，编码器只处理可见块，轻量解码器再结合掩码标记预测缺失块的像素；原论文的均方误差只计算在被遮蔽块上\scite{ka-he2022mae}
因此，编码器没有看到的位置与损失评价的位置恰好互补。MAE也研究过按块归一化后的像素目标，但无论是否归一化，答案都可在训练前从原图确定，仍属于观测目标。

遮蔽比例控制信息与难度的平衡。比例太低时，模型可能只做局部插值；比例太高时，一个给定输入可能对应许多合理恢复结果，点预测会平均这些可能性。随机遮蔽、连续片段遮蔽和结构化遮蔽还改变了可利用的上下文。具体比例不能仅按训练损失选择：更容易恢复通常会降低损失，却未必更能保留下游任务所需的信息。

构造过程还要检查预处理顺序。若先用覆盖目标位置的平滑、插值或全图归一化统计产生输入，再把目标位置标成掩码，答案可能已经通过邻居或统计量泄露。安全的原则是明确每个预处理算子的感受范围，在拆分输入与目标后重新审查它是否读取了$M$中的内容。

如果掩码位置改为预测另一个编码器生成的特征，就不能仅凭“遮蔽”认定目标来自观测。应继续追踪特征编码器是否固定、是否随训练更新，以及当前损失是否向它回传梯度；这类动态表示目标属于第~\ref{sec:data-derived-knowledge-training-objectives}节的分析对象。

\subsubsection{基于扰动的去噪恢复}
\label{sec:data-derived-knowledge-denoising-recovery}

\term{去噪恢复}（denoising reconstruction）从扰动分布$q(\widetilde{\bm{x}}\mid\bm{x})$抽取受损输入$\widetilde{\bm{x}}$，让模型恢复扰动前保留的$\bm{x}$。这里的“噪声”可以是数值加性噪声、词元替换、局部打乱、颜色缺失或其他人为损坏，不等同于数据采集过程中的物理噪声。

若$R\subseteq I$是已知受扰位置，一个一般形式为
\begin{equation}
  \mathcal{L}_{\mathrm{denoise}}
  =\mathbb{E}_{\widetilde{\bm{x}}\sim q(\cdot\mid\bm{x})}
   \left[\frac{1}{|R|}\sum_{j\in R}\ell\!\left(f_{\bm{\theta}}(\widetilde{\bm{x}})_j,x_j\right)\right].
  \label{eq:data-derived-knowledge-denoising}
\end{equation}
改为对$I$求和时，目标要求模型同时保持未扰部分和恢复受扰部分。两种范围会给出不同梯度，报告结果时不能只写“使用去噪损失”。

扰动太弱时，最优策略可能接近恒等复制；扰动太强时，$\widetilde{\bm{x}}$不足以确定$\bm{x}$，模型只能按损失选择条件均值、众数或分布。后者不是优化失败，而是输入信息本身不充分。例如一块花瓣完全被随机纹理替换后，颜色可能有多个合理答案；均方误差倾向输出这些答案的平均，生成分布则可以保留多种可能。

在线生成扰动只需保存原样本和随机配置，却增加每轮数据处理；离线保存多个损坏副本节省生成时间，却增加存储并限制扰动多样性。无论采用哪种方式，都应记录随机种子或分布参数，并让评价扰动与训练扰动的关系可追溯。只在同一种教学噪声上恢复良好，不能证明模型能处理真实采集误差。

还要分开评价恢复质量与表示效用。模型可能准确撤销某种固定压缩伪影，却没有保留花卉类别所需的形态；也可能像MAE那样对像素逐点恢复并不唯一，却学到对识别有用的编码。训练目标提供学习依据，下游用途是否受益仍需独立证据。

\subsection{基于上下文预测的目标构造}
\label{sec:data-derived-knowledge-context-prediction}

有些记录天然带有顺序或邻接结构。此时无需随机选择任意掩码，可以约定一个位置的上下文，再把相邻片段作为目标。\term{上下文预测}就是根据给定前缀、邻域或其他可见内容预测记录中的目标片段。它与掩码恢复的区别不在于都预测词元，而在于可见方向和切分规则。

对词元序列$x_1,\ldots,x_T$，只允许第$t$个预测读取前缀$x_{<t}$时，逐位置负对数似然为
\begin{equation}
  \mathcal{L}_{\mathrm{context}}
  =-\sum_{t=1}^{T}\log p_{\bm{\theta}}(x_t\mid x_{<t}).
  \label{eq:data-derived-knowledge-context-prediction}
\end{equation}
这一分解与第~\ref{sec:pgm-static-directed-autoregressive}节的自回归模型一致；本节关注的是怎样从一条记录切出输入和答案。训练记录“花朵逐渐开放”可以形成“花朵$\to$逐渐”和“花朵逐渐$\to$开放”等多个位置目标，监督覆盖随可预测位置增多，但计算也近似按位置数增长。

遮蔽词元时，模型通常可读取目标左右两侧；前缀预测则只能读取目标之前的内容。BERT的掩码词元目标属于前一种安排：原词元从记录中保留，输入使用双向上下文\scite{devlin2019}
二者即使都采用交叉熵，也不能互换条件概率的含义。

训练时常把真实前缀$x_{<t}$交给模型，这称为教师强制式输入安排；连续生成时，后续位置读到的可能是模型自己此前的输出。于是，式~\eqref{eq:data-derived-knowledge-context-prediction}上的一步预测误差不能直接代表长序列生成误差。窗口截断还会让模型看不到更早信息，位置选择和片段边界也可能制造不自然的上下文。

上下文不局限于文本。图像可以根据左侧区域预测右侧区域，音频可以根据相邻频段预测被保留片段，多模态记录也可以让一个模态预测同一事件的另一观测。判断是否仍属于观测目标，要看答案能否从已记录数据中直接取出，而不是看输入和输出是否属于同一模态。

常见捷径来自局部重复和预处理泄露。网页模板、固定边框或传感器周期背景可能让模型不理解主体也能预测相邻内容；分词器、缓存特征或裁剪坐标若编码了目标，也会夸大训练效果。改变上下文范围、删除可疑元数据并在独立下游任务检验，能够区分“会利用记录规律”和“获得了预期知识”。

\subsection{基于时序预测的目标构造}
\label{sec:data-derived-knowledge-temporal-prediction}

上下文顺序可以只是文本排列约定，时序预测则额外要求每个输入都对应一个真实判断时点。\term{时序预测}（temporal prediction）用时点$t$以前的历史观测预测$t$之后的记录。未来部分在整理训练数据时已经存在，因此仍是观测目标；模型部署到时点$t$时，这些内容尚未发生或尚未观测。

设历史窗口长度为$L$、起始预测跨度为$h\geq1$、目标轨迹长度为$H$。输入和目标分别写为
\begin{equation}
  \bm{x}_{t-L+1:t}
  \quad\longrightarrow\quad
  \bm{y}_{t,h,H}
  =\bm{x}_{t+h:t+h+H-1}.
  \label{eq:data-derived-knowledge-temporal-windows}
\end{equation}
必须同时给出采样间隔，因为“未来十步”在每分钟采样和每天采样的记录中含义不同。预测一个时点时$H=1$；预测整段植物生长轨迹时$H>1$，输出维数和误差相关性都随之改变。

未来结果近似唯一且平方误差符合使用代价时，可以做点预测。若同一历史后可能出现多条合理轨迹，条件分布、分位数或离散区间目标更能表达不确定性。选择分布目标并不会消除不可预测性，只是避免把所有可能性压成一个均值；评价仍要与实际决策需要对应。

短跨度常依赖局部连续性，长跨度则更容易受未观测干预、累积误差和环境变化影响。比较不同跨度时，应固定数据范围并分别报告，不能把多个$h$的误差平均后掩盖远期失效。多个跨度共同训练会增加监督覆盖，也会增加输出和存储成本。

时序任务最危险的错误是从未来偷取信息。窗口随机打散后再按样本切分，可能让同一轨迹的高度重叠窗口同时进入训练和测试；对全序列计算归一化均值、用未来点插值当前缺失值，都会把判断时点后的信息带回输入。应先按未来使用单位划分时间、植株、设备或轨迹，再只用训练时可获得的统计量拟合预处理。

预测收益也不能等同于知识质量。重复背景可以让下一帧像素误差很低，却不表示模型理解植物状态；相反，长期未来本就具有较大随机性，误差较高也不一定说明表示毫无价值。训练目标的可预测程度、模型学到的信息和下游用途需要分开评价。

\section{基于样本关系的知识获取}
\label{sec:data-derived-knowledge-relations}

观测内容提供的是“这个位置原来是什么”。另一类依据不直接给出内容答案，而是说明两个观测之间应保持什么关系。\term{关系知识}就是可追溯的样本关系记录，包括应当保留共同信息的正对应，以及在相对判别中作为参照的负对照；它们可以来自程序变换、采集记录或采样规则。关系知识可以在编码前保存，而表示目标要等模型完成当前步前向计算后才能得到。

关系本身与使用关系的目标函数是两个层次。同一张花卉照片的两个裁剪构成正对应后，可以用SimCLR在候选集合中识别配对，也可以用BYOL预测另一视图的表示，还可以用VICReg对称匹配并约束批内统计。反过来，看到一个余弦损失也不能推断正对应是怎样建立的。

本节沿同一批原样本追踪三个对象：可提前保存的关系记录、当前训练步由模型计算的目标值、跨训练步变化的参数。关系假设是否合理、当前计算图是否落实关系、训练后表示是否有用，分别需要构造证据、梯度检查和下游评价，不能由一次损失下降同时回答。

\subsection{样本关系的构造}
\label{sec:data-derived-knowledge-relation-construction}

一条可追溯关系至少应记录原对象、关系两侧、配对或排除依据、关系适用范围和构造配置。例如“同一原图的两个裁剪”只说明来源一致，还要记录裁剪是否都覆盖主体；“同一时刻的音频和画面”还要记录允许的同步误差。无法确定的关系可以保留置信范围，不必强行变成二值标签。

图~\ref{fig:data-derived-knowledge-relation-construction}区分数据生成、正对应与需要检查的关系。蓝色箭头由原记录生成视图，青色实线表示希望保留共同信息的正对应；左侧灰色虚线是负对照假设，右侧赭色虚线则是教学设定中的错配记录。程序生成并不证明裁剪保持了任务语义，自然配对也不证明图文内容逐项对应。

\begin{figure*}[htbp]
  \centering
  \begin{tikzpicture}[
  x=1mm,y=1mm,
  font=\figurefont\footnotesize,
  panel/.style={draw=BookRule,fill=NNPanel,line width=.85pt,rounded corners=1.5mm},
  item/.style={draw=FigureInput,fill=FigureInput!10,line width=1.1pt,minimum width=24mm,minimum height=11mm,align=center,inner sep=1.5mm},
  textitem/.style={draw=FigureOutput,fill=FigureOutput!8,line width=1.1pt,minimum width=28mm,minimum height=11mm,align=center,inner sep=1.5mm},
  generate/.style={knowledge arrow,draw=FigureInput,line width=1.1pt},
  positive/.style={draw=FigureModel,line width=1.4pt},
  negative/.style={draw=BookMuted,line width=.9pt,dashed},
  warning/.style={font=\figurefont\scriptsize,text=FigureLoss,align=center}
]
  \draw[panel] (1,-46) rectangle (82,10);
  \node[font=\figurefont\bfseries\footnotesize,text=BookInk] at (41.5,5) {(a) 数据变换建立正对应};
  \node[item] (original) at (12.5,-13) {原照片\\花朵与叶片};
  \node[item] (crop-a) at (45,-5) {裁剪A\\花朵};
  \node[item] (crop-b) at (45,-22) {裁剪B\\花朵与叶片};
  \node[item,draw=BookAmber,fill=BookAmber!10] (same-class) at (68,-35) {同类另一株\\未必是负语义};
  \draw[generate] (original) -- (crop-a);
  \draw[generate] (original) -- (crop-b);
  \draw[positive] (crop-a.south) -- node[right,text=FigureModel] {正对应：同一原对象} (crop-b.north);
  \draw[negative] (crop-b.south east) -- node[above,sloped,text=BookMuted] {候选负对照} (same-class.north west);
  \node[warning] at (24,-38) {关系可能失配：裁剪只剩背景，\\或同类样本被当作负例。};

  \draw[panel] (87,-46) rectangle (167.5,10);
  \node[font=\figurefont\bfseries\footnotesize,text=BookInk] at (127.25,5) {(b) 采集记录建立自然配对};
  \node[item] (photo) at (101,-13) {花卉照片\\白色花瓣};
  \node[textitem] (caption) at (136,-5) {原有描述\\“白色花瓣，长叶”};
  \node[textitem,draw=BookAmber,fill=BookAmber!10] (wrong-caption) at (136,-24) {错配描述\\“黄色花瓣”};
  \node[item] (other-photo) at (154.5,-37) {其他照片\\负对照候选};
  \draw[positive] (photo) -- node[above,sloped,text=FigureModel] {正对应：记录配对} (caption);
  \draw[negative] (photo) -- node[below,sloped,text=BookMuted] {错误关系} (wrong-caption);
  \draw[negative] (caption.south east) -- node[right,text=BookMuted] {候选负对照} (other-photo.north);
  \node[warning,anchor=west] at (89,-38) {配对只保证采集关联，\\不保证所有内容逐项对应。};

  \draw[generate] (3,-52) -- (19,-52);
  \node[anchor=west,text=BookMuted] at (21,-52) {数据生成};
  \draw[positive] (55,-52) -- (71,-52);
  \node[anchor=west,text=BookMuted] at (73,-52) {拟保留共同信息};
  \draw[negative] (117,-52) -- (133,-52);
  \node[anchor=west,text=BookMuted] at (135,-52) {负对照假设};
\end{tikzpicture}

  \caption{正对应与负对照的构造依据。左侧从同一照片生成两种视图，并接入同类另一株作为负对照候选；右侧对照原有图文配对与一条错配描述。负对照是训练采用的关系假设，不表示已经确认不同语义；错配记录则不能用作正对应。}
  \label{fig:data-derived-knowledge-relation-construction}
\end{figure*}
\FloatBarrier

\subsubsection{基于数据变换的正对应关系构造}
\label{sec:data-derived-knowledge-transformation-positive}

从同一原样本抽取两次或多次随机变换，可以低成本地产生正对应。构造假设是：变换改变了不希望表示敏感的因素，同时保留了希望共享的信息。整体视图、局部裁剪、尺度变化、颜色扰动、语音加噪和文本改写各自改变的内容不同，不能只因都由程序生成就使用同一强度。

以花卉照片为例，两次裁剪都包含花朵时，正对应可以鼓励表示忽略取景差异；一路只剩背景时，两路未必有足够共同内容；强颜色变化若抹去了品种判断所需的色彩，则“不变性”与下游任务冲突。若下游任务需要方向或位置，也不能无条件使用旋转、翻转或大范围平移。

因此，变换配置不是普通的数据清洗参数，而是对“哪些变化不应改变表示”的知识声明。应分别记录变换种类、组合顺序、强度、随机参数和主体覆盖情况，并用抽样可视化或自动规则检查正对应是否仍有共同内容。视图数量增加会提供更多对应，却也按编码次数增加计算。

关系可以一对多。一个整体视图可能与多个局部视图对应，但每个局部视图只覆盖整体的一部分。后续目标若要求局部表示完整复制整体信息，就超出了关系本身支持的范围；DINO的局部到全局匹配便需要依赖多裁剪设计和训练经验，而不是由“同一原图”逻辑保证成功。

本小节只建立正对应，不规定是否使用负对照、教师分支或批内统计。具体的目标形式由第~\ref{sec:data-derived-knowledge-training-objectives}节比较。

\subsubsection{基于自然配对的正对应关系构造}
\label{sec:data-derived-knowledge-natural-positive}

\term{自然配对}是采集过程中已经存在的对应记录，例如同一对象的多次拍摄、同步音视频、图片及说明文字、相邻时刻的传感器读数。与数据变换相比，它不要求两侧属于同一数据类型，却更依赖对象标识、时间戳、事件编号或文档结构是否可靠。

花卉照片与说明文字可能有三种范围：描述整株形态，只描述花色，或者误配到另一张照片。前两种都是真实配对，却只在不同内容上对应；第三种则是关系错误。训练前应说明关系是整例对应、局部对应还是允许一个对象匹配多条记录，并保留缺失一侧和多重匹配。

% 完整作者和会议条目不能跨页续排。
\Needspace{65mm}
CLIP使用图像与文本的已有配对构造正对应，再把同一批次中的其他图文组合作为对照\scite{radford2021}
文字给出了人类书写的描述信息，因此这类学习虽然不要求为目标分类逐图标注，却不能笼统称为“没有人类监督”。本章关心的是配对如何成为训练依据，而不是把所有非类别标签信号归入同一成本口径。

自然配对的常见误差包括时间同步偏移、描述不完整、重复记录和共享背景。若多张照片复用同一句模板描述，把它们互相当作负对照会制造冲突；若音频比画面延迟，固定时刻配对会把不同事件连在一起。扩大数据量不能自动修复系统性错配，需要回查采集协议和关联记录。

多模态输入还带来额外编码与存储成本。应分别报告两侧的缺失率、配对核验方式和模态处理开销，避免只按“无需新标签”评价其投入。

\subsubsection{基于样本采样的负对照关系构造}
\label{sec:data-derived-knowledge-negative-sampling}

对比判别需要回答“正确对应相对于哪些候选更匹配”。为此，可以从小批量、缓存或检索集合选择\term{负对照}。负对照表示当前目标允许它与查询相对区分，不等于已经获得“语义不同”的人工标签。

随机抽取通常覆盖广，但很多候选与查询差异过大，产生的梯度较弱。根据当前表示检索相似候选可以得到\term{难负例}，让目标关注细微差别；然而，表示相似只说明模型当前难以区分，并不证明它们语义不同。反复检索少数候选还会让训练过度适应当前表示的偏差。

同一植株的另一张照片、同品种不同植株和完全不同对象，在实例区分任务中可以承担不同角色。若目标是识别具体植株，前两者的关系不同；若目标是识别品种，把同品种样本作为负对照会冲突。关系集合应先由任务需要确定，不能把“不是已知正例”自动改写成“确定负例”。

批内对比没有单独的采样文件，仍然执行了采样：小批量组成决定了分母中的候选。重复记录、类别失衡和同源样本会改变负对照分布；缓存还会保存由旧参数计算的过时表示。扩大候选集合能增加区分压力，也会增加错误负例、通信和存储成本。

一个查询存在多个合理正例时，应把它们显式加入正对应集合，或采用允许多正例的目标。否则，其余合理对应进入分母后会受到排斥。负对照是关系构造结果，对比损失是使用结果的方式；两者必须分别检查。

\subsection{训练目标的构造}
\label{sec:data-derived-knowledge-training-objectives}

关系确定后，还要回答三个问题：模型比较什么量，哪一路产生目标，哪些参数接收当前损失的梯度。比较量可以是相对相似度、另一视图的输出，也可以包括一批表示的统计量。下面分别建立这些计算，再比较训练结束后应保留哪个表示。交换两视图另算一项损失，与同一项损失向两路回传，是两个不同问题。

% 标题、表示定义及其边注在同一页内保留足够高度。
\Needspace{80mm}
\subsubsection{基于对比判别的目标构造}
\label{sec:data-derived-knowledge-contrastive}

\term{对比判别}让正对应相对于候选集合获得更高相似度。以SimCLR为例，同一批$n$个原样本各生成两个视图，得到$2n$个编码。编码器$f_{\bm{\theta}}$输出下游表示$\bm{h}_i$，投影器$g_{\bm{\theta}}$再给出训练损失使用的$\bm{z}_i=g_{\bm{\theta}}(\bm{h}_i)$。训练结束后通常使用$\bm{h}_i$，不能把投影空间与下游表示混成一个对象\scite{chen2020}

以下假定$n\geq2$，且参与比较的投影表示均非零。定义余弦相似度
\[
  \operatorname{sim}(\bm{z}_i,\bm{z}_k)
  =\frac{\langle\bm{z}_i,\bm{z}_k\rangle}
  {\lVert\bm{z}_i\rVert_2\lVert\bm{z}_k\rVert_2}.
\]
若$j$是锚点$i$的唯一正对应，SimCLR的单向归一化温度交叉熵损失为
\begin{equation}
  \ell_{i,j}
  =-\log
  \frac{\exp\!\left(\operatorname{sim}(\bm{z}_i,\bm{z}_j)/\tau\right)}
  {\displaystyle\sum_{k=1}^{2n}\mathbb{I}[k\neq i]\exp\!\left(\operatorname{sim}(\bm{z}_i,\bm{z}_k)/\tau\right)},
  \qquad \tau>0.
  \label{eq:data-derived-knowledge-simclr}
\end{equation}
分母排除锚点自身，却包含正对应$j$和其余$2n-2$个候选。完整目标让两个视图轮流作锚点，对全部$2n$项求平均。较小温度会更强调相似度差异，但也放大难负例和错误关系的影响。

一个$n=2$的教学批次含$A_1,A_2,B_1,B_2$四个视图，正对应为$(A_1,A_2)$与$(B_1,B_2)$。设余弦相似度矩阵为
\[
  \bm{S}=
  \begin{bmatrix}
  1&0.8&0.1&0\\
  0.8&1&0&-0.1\\
  0.1&0&1&0.7\\
  0&-0.1&0.7&1
  \end{bmatrix}.
\]
取$\tau=0.5$，以$A_1$为锚点时，正确配对概率为$\exp(1.6)/[\exp(1.6)+\exp(0.2)+\exp(0)]\approx0.690$，所以$\ell_{A_1,A_2}\approx0.371$。这个概率只表示在当前候选集合中支持正确配对的相对程度，不是“属于同一花卉类别”的概率。

批内其他视图同时提供对照，使损失依赖小批量组成。扩大批量通常增加候选与相似度计算，分布式训练还需交换表示；小批量若缺少有区分力的候选，目标可能变弱。缓存和检索可以改变候选规模，却带来旧表示与采样偏差。

错误负例会直接进入式~\eqref{eq:data-derived-knowledge-simclr}的分母。同品种花卉若被当作负对照，实例区分目标可能主动排斥下游希望聚合的样本。多正例目标、去重和关系过滤可以减轻冲突，但都需要额外关系信息，不能由损失自行判断。

若所有投影表示都等于同一个非零向量$\bm{u}$，全部余弦相似度都是$1$，于是
\[
  \ell_{i,j}=\log(2n-1).
\]
这个常数输出没有配对优势。只要表示空间和模型允许正对应比负对照更相似，就能取得更小损失，例如上面的教学批次有$\ell_{A_1,A_2}\approx0.371<\log3\approx1.099$。但目标值较差不等于优化会自动离开该点：当两向量相同时，余弦相似度对任一非零输入的梯度为零，因此精确的常数输出仍是损失关于表示的驻点。若$n=1$，分母中只剩正对应，损失更是恒为零。对比目标提供的是有条件的区分偏好，不是任意初始化下避免坍塌的保证；即使学出了差异，也仍要检查它来自主体还是背景、设备或水印。

\subsubsection{基于跨视图预测的目标构造}
\label{sec:data-derived-knowledge-cross-view-prediction}

\term{跨视图预测}只利用正对应，让一个视图的在线分支预测另一个视图产生的表示目标。BYOL把在线分支分成编码器$f_{\bm{\theta}}$、投影器$g_{\bm{\theta}}$和预测器$q_{\bm{\theta}}$，目标分支只有参数为$\bm{\xi}$的编码器和投影器。对视图$\bm{v}_A,\bm{v}_B$，记
\[
  \bm{p}_A=q_{\bm{\theta}}\!\left(g_{\bm{\theta}}(f_{\bm{\theta}}(\bm{v}_A))\right),
  \qquad
  \bm{z}^{\mathrm{tar}}_B=g_{\bm{\xi}}(f_{\bm{\xi}}(\bm{v}_B)).
\]
单向损失比较两个$\ell_2$归一化向量：
\begin{equation}
  \ell_{A\to B}
  =\left\lVert
  \frac{\bm{p}_A}{\lVert\bm{p}_A\rVert_2}
  -\operatorname{sg}\!\left(
  \frac{\bm{z}^{\mathrm{tar}}_B}{\lVert\bm{z}^{\mathrm{tar}}_B\rVert_2}
  \right)\right\rVert_2^2
  =2-2\frac{\langle\bm{p}_A,\bm{z}^{\mathrm{tar}}_B\rangle}
  {\lVert\bm{p}_A\rVert_2\lVert\bm{z}^{\mathrm{tar}}_B\rVert_2}.
  \label{eq:data-derived-knowledge-byol}
\end{equation}
$\operatorname{sg}$表示停止梯度：当前反向传播把目标向量当作常量。BYOL还交换视图计算$\ell_{B\to A}$，总损失为两项之和；但每一项都只沿相应在线分支求导。对称化的求和没有消除在线与目标分支的更新差别\scite{ka-grill2020byol}

在线参数完成梯度更新后，目标编码器和投影器的对应参数用指数移动平均更新：
\begin{equation}
  \bm{\xi}^{(t+1)}
  =m\bm{\xi}^{(t)}+(1-m)\bm{\theta}_{\mathrm{enc+proj}}^{(t+1)},
  \qquad 0\leq m\leq1.
  \label{eq:data-derived-knowledge-byol-ema}
\end{equation}
预测器没有对应的目标分支，其参数只由梯度更新。停止梯度规定当前步的求导路径，EMA规定下一步目标参数怎样变化；两者不是同一个操作。

这里预测的是学得向量，不是目标视图的像素。像素在编码前就能从记录中保留，属于第~\ref{sec:data-derived-knowledge-observation}节；$\bm{z}^{\mathrm{tar}}_B$则由当前目标网络计算，会随训练改变。它也不同于带预设任务含义的分类教师输出：BYOL目标向量的坐标没有预先命名为花卉类别。

BYOL预训练结束后通常保留在线编码器$f_{\bm{\theta}}$，将其输出作为下游特征；投影器、预测器和目标分支都不进入这一下游接口。因而训练时匹配的$\bm{p}_A$与$\bm{z}^{\mathrm{tar}}_B$不是最终交给分类器的表示。

DINO提供了另一种跨视图表示目标。学生和教师都由骨干编码器及投影头组成，结构相同、参数不同，但没有BYOL式预测头。设投影头输出$K$维向量$\bm{u}_{\mathrm{s}}(\bm{v})$和$\bm{u}_{\mathrm{t}}(\bm{v})$，教师中心为$\bm{c}$，则学生与教师分布分别为
\begin{align}
  p_{\mathrm{s},k}(\bm{v})
  &=\frac{\exp(u_{\mathrm{s},k}(\bm{v})/\tau_{\mathrm{s}})}
  {\sum_{r=1}^{K}\exp(u_{\mathrm{s},r}(\bm{v})/\tau_{\mathrm{s}})},\\
  p_{\mathrm{t},k}(\bm{v})
  &=\frac{\exp((u_{\mathrm{t},k}(\bm{v})-c_k)/\tau_{\mathrm{t}})}
  {\sum_{r=1}^{K}\exp((u_{\mathrm{t},r}(\bm{v})-c_r)/\tau_{\mathrm{t}})}.
  \label{eq:data-derived-knowledge-dino-distributions}
\end{align}
教师只处理两个全局视图，学生处理全局与局部视图；每个教师全局视图与除同一视图外的学生输出计算交叉熵。教师端停止梯度并用学生参数的EMA更新，中心$\bm{c}$也用教师批量输出的移动平均更新\scite{ka-caron2021dino}

DINO的居中抑制某个输出维度长期支配，较低教师温度进行锐化，防止目标总是接近均匀分布。原论文把两者与动量教师、多裁剪共同使用，并通过消融观察其作用；不能把“居中”或“低温”单独写成对任意模型都成立的防坍塌定理。$K$维输出也没有人工类别名称，低交叉熵不表示自动发现了$K$个真实物种。

DINO下游使用的是骨干特征，投影头及其$K$维分布用于预训练。本章的下游接口约定选用训练结束时的教师编码器检查点，并舍弃投影头；若改用学生检查点，也应明确记录分支，不能把两个编码器的结果混在同一评价中。

跨视图预测省去显式负对照，却增加目标分支或多视图计算。若两个裁剪共同信息不足，学生只能拟合目标分支中可预测的部分；若表示全部变成常量，单纯的匹配项仍可能很小。完整优化动态怎样控制这种退化，要结合第~\ref{sec:data-derived-knowledge-representation-degeneration}节分析。

\subsubsection{基于对称匹配的目标构造}
\label{sec:data-derived-knowledge-symmetric-matching}

\term{对称匹配}在同一个配对项中平等对待两路表示：减小$\bm{z}_i$与$\bm{z}'_i$的差异时，梯度同时回到两路。两路可以共享编码器参数，也可以使用不同结构；不需要先指定一个独立教师产生答案，也不必停止某一路梯度。这里的“对称”描述单项损失的求导角色，不只是把两个方向的非对称损失相加。

VICReg先用编码器$f_{\bm{\theta}}$取得特征，再用\term{扩展头}（expander）$h_{\bm{\phi}}$映射到计算训练损失的嵌入空间，即$\bm{z}_i=h_{\bm{\phi}}(f_{\bm{\theta}}(\bm{v}_i))$。这里采用两路共享编码器与扩展头的安排，$\bm{\theta}$和$\bm{\phi}$都接收两路损失的梯度。将两个正对应视图批次的嵌入按行排列，得到$\bm{Z},\bm{Z}'\in\mathbb{R}^{n\times d}$，并要求$n\geq2$。匹配项为
\begin{equation}
  s(\bm{Z},\bm{Z}')
  =\frac{1}{n}\sum_{i=1}^{n}\lVert\bm{z}_i-\bm{z}'_i\rVert_2^2.
  \label{eq:data-derived-knowledge-vicreg-invariance}
\end{equation}
若只有这一项，所有样本输出同一常量就能取得零损失。VICReg因此分别在两路加入方差与协方差约束\scite{ka-bardes2022vicreg}

对第$j$维，本章的批内方差采用样本方差分母$n-1$。令$\bar z_j=n^{-1}\sum_i z_{i,j}$，带数值稳定项的标准差为
\[
  \begin{aligned}
    \operatorname{Var}(\bm{Z}_{:,j})
    &=\frac{1}{n-1}\sum_{i=1}^{n}(z_{i,j}-\bar z_j)^2,\\
    \sigma_j(\bm{Z})
    &=\sqrt{\operatorname{Var}(\bm{Z}_{:,j})+\epsilon}.
  \end{aligned}
\]
取$\epsilon>0$、$\gamma>\sqrt{\epsilon}$，使常数批次的标准差仍低于目标阈值。方差惩罚为
\begin{equation}
  v(\bm{Z})
  =\frac{1}{d}\sum_{j=1}^{d}\max(0,\gamma-\sigma_j(\bm{Z})).
  \label{eq:data-derived-knowledge-vicreg-variance}
\end{equation}
它只在标准差低于阈值$\gamma$时施加惩罚，是与其他项共同优化的软约束，并不保证每一步都满足$\sigma_j\geq\gamma$。

令$\overline{\bm{z}}=n^{-1}\sum_i\bm{z}_i$，样本协方差矩阵为
\[
  \bm{C}(\bm{Z})
  =\frac{1}{n-1}\sum_{i=1}^{n}
  (\bm{z}_i-\overline{\bm{z}})
  (\bm{z}_i-\overline{\bm{z}})^{\top}.
\]
协方差惩罚只处理非对角元素：
\begin{equation}
  c(\bm{Z})
  =\frac{1}{d}\sum_{\substack{j,k=1\\j\neq k}}^{d}C_{j,k}(\bm{Z})^2.
  \label{eq:data-derived-knowledge-vicreg-covariance}
\end{equation}
完整损失为
\begin{equation}
  \mathcal{L}_{\mathrm{VICReg}}
  =\lambda s(\bm{Z},\bm{Z}')
  +\mu\!\left[v(\bm{Z})+v(\bm{Z}')\right]
  +\nu\!\left[c(\bm{Z})+c(\bm{Z}')\right],
  \label{eq:data-derived-knowledge-vicreg}
\end{equation}
其中$\lambda,\mu,\nu>0$控制三类要求。匹配项约束逐样本正对应，方差与协方差项则跨样本估计每一路的分布性质；后两者不是把其他样本逐对标成负例。

令$\bm{Z}=\bm{Z}'$且所有行都相同，匹配项与协方差项为零，而每路的方差惩罚为$\max(0,\gamma-\sqrt{\epsilon})$。在上述参数条件下它为正；取教学数值$\gamma=1$、$\epsilon=10^{-4}$，便有$v(\bm{Z})=0.99$。若不加这一条件，例如$\gamma=\epsilon=1$，常数批次的三项损失可以全部为零。即使惩罚为正，精确常数批次的方差梯度仍为零，不能由此推出优化必能逃离该驻点。若各行沿直线$(a,a)$变化，两维都有方差，但协方差非零，说明坐标重复携带同一变化。

批量太小时，方差和协方差估计波动较大；维数$d$较大时，显式协方差计算和存储也有成本。即使所有维度都保留方差且线性相关较低，表示仍可能主要编码背景、设备或其他无关因素。统计惩罚提高若干退化状态的损失，不提供下游语义或优化收敛保证。VICReg预训练后舍弃扩展头$h_{\bm{\phi}}$，保留编码器$f_{\bm{\theta}}$的输出供下游使用。

至此可以把三类目标放回同一计算图比较。图~\ref{fig:data-derived-knowledge-relation-objectives}按前向值、当前步梯度和跨步参数平均分开布线：BYOL的目标分支参与前向计算，但不接收当前损失的梯度；VICReg则让同一匹配项向两路回传。交换视图求和不能替代这一区分。

\begin{figure*}[htbp]
  \centering
  \begin{tikzpicture}[
  x=1mm,y=1mm,
  font=\figurefont\footnotesize,
  panel/.style={draw=BookRule,fill=NNPanel,line width=.85pt,rounded corners=1.5mm},
  input/.style={draw=NNInput,fill=NNInput!16,line width=1pt,minimum width=7mm,minimum height=7mm,inner sep=1mm,align=center},
  compute/.style={draw=NNHidden,fill=NNHidden!16,line width=1pt,minimum width=10mm,minimum height=8mm,inner sep=1mm,align=center,rounded corners=1mm},
  output/.style={draw=NNOutput,fill=NNOutput!16,line width=1pt,minimum width=8mm,minimum height=8mm,inner sep=1mm,align=center},
  loss/.style={draw=NNGradient,fill=NNGradient!12,line width=1pt,minimum width=10mm,minimum height=8mm,inner sep=1mm,align=center,rounded corners=1mm},
  forward/.style={knowledge arrow,draw=NNWire,line width=1pt},
  gradient/.style={knowledge arrow,draw=NNGradient,line width=1.1pt,dashed,rounded corners=1mm},
  gradient stop/.style={-{Bar[width=2.4mm]},draw=NNGradient,line width=1.1pt,dashed,rounded corners=1mm},
  ema/.style={knowledge arrow,draw=BookTeal,line width=1pt,dash dot},
  note/.style={font=\figurefont\scriptsize,text=BookMuted,align=center}
]
  \draw[panel] (.5,103) rectangle (168.5,153);
  \draw[panel] (.5,48) rectangle (168.5,98);
  \draw[panel] (.5,-10) rectangle (168.5,42);
  \node[font=\figurefont\bfseries\footnotesize,text=BookInk,anchor=west] at (5,148) {(a) 对比判别：SimCLR};
  \node[note,anchor=east] at (163,148) {两路共享编码器与投影器};
  \node[input] (s-v1) at (10,132) {$\bm{v}_1$};
  \node[input] (s-v2) at (10,117) {$\bm{v}_2$};
  \node[compute,minimum width=20mm] (s-f1) at (42,132) {$f_{\bm{\theta}}$};
  \node[compute,minimum width=20mm] (s-f2) at (42,117) {$f_{\bm{\theta}}$};
  \node[compute,minimum width=20mm] (s-g1) at (76,132) {$g_{\bm{\theta}}$};
  \node[compute,minimum width=20mm] (s-g2) at (76,117) {$g_{\bm{\theta}}$};
  \node[output,minimum width=14mm] (s-z1) at (108,132) {$\bm{z}_1$};
  \node[output,minimum width=14mm] (s-z2) at (108,117) {$\bm{z}_2$};
  \node[loss,minimum width=26mm] (s-loss) at (148,124.5) {相似度$\bm{S}$\\与配对损失};
  \draw[forward] (s-v1.east)--(s-f1.west);
  \draw[forward] (s-f1.east)--(s-g1.west);
  \draw[forward] (s-g1.east)--(s-z1.west);
  \draw[forward] (s-v2.east)--(s-f2.west);
  \draw[forward] (s-f2.east)--(s-g2.west);
  \draw[forward] (s-g2.east)--(s-z2.west);
  \draw[NNWire,line width=1pt,rounded corners=1mm] (s-z1.east) -- (126,132) -- (126,124.5);
  \draw[NNWire,line width=1pt,rounded corners=1mm] (s-z2.east) -- (126,117) -- (126,124.5);
  \draw[forward] (126,124.5)--(s-loss.west);
  \draw[gradient] (s-loss.north) -- (148,140) -- (42,140) -- (s-f1.north);
  \draw[gradient] (76,140)--(s-g1.north);
  \draw[gradient] (108,140)--(s-z1.north);
  \draw[gradient] (s-loss.south) -- (148,109) -- (42,109) -- (s-f2.south);
  \draw[gradient] (76,109)--(s-g2.south);
  \draw[gradient] (108,109)--(s-z2.south);

  \node[font=\figurefont\bfseries\footnotesize,text=BookInk,anchor=west] at (5,93) {(b) 跨视图预测：BYOL};
  \node[note,anchor=east] at (163,93) {只画$A\to B$；交换视图另算一项};
  \node[input] (b-va) at (10,78) {$\bm{v}_A$};
  \node[compute,minimum width=20mm] (b-fo) at (42,78) {$f_{\bm{\theta}}$};
  \node[compute,minimum width=20mm] (b-go) at (76,78) {$g_{\bm{\theta}}$};
  \node[compute,minimum width=20mm] (b-qo) at (108,78) {$q_{\bm{\theta}}$};
  \node[input] (b-vb) at (10,59) {$\bm{v}_B$};
  \node[compute,minimum width=20mm] (b-ft) at (42,59) {$f_{\bm{\xi}}$};
  \node[compute,minimum width=20mm] (b-gt) at (76,59) {$g_{\bm{\xi}}$};
  \node[output,minimum width=20mm] (b-sg) at (108,59) {$\operatorname{sg}$};
  \node[loss,minimum width=20mm] (b-loss) at (148,68.5) {$\ell_{A\to B}$};
  \draw[forward] (b-va.east)--(b-fo.west);
  \draw[forward] (b-fo.east)--(b-go.west);
  \draw[forward] (b-go.east)--(b-qo.west);
  \draw[forward,rounded corners=1mm] (b-qo.east) -| (b-loss.north);
  \draw[forward] (b-vb.east)--(b-ft.west);
  \draw[forward] (b-ft.east)--(b-gt.west);
  \draw[forward] (b-gt.east)--(b-sg.west);
  \draw[forward,rounded corners=1mm] (b-sg.east) -| (b-loss.south);
  \draw[gradient] (b-loss.east) -- (164,68.5) -- (164,86) -- (42,86) -- (b-fo.north);
  \draw[gradient] (76,86)--(b-go.north);
  \draw[gradient] (108,86)--(b-qo.north);
  \draw[gradient stop] (164,68.5) -- (164,52) -- (108,52) -- (b-sg.south);
  \draw[ema] (b-fo.south)--(b-ft.north);
  \draw[ema] (b-go.south)--(b-gt.north);
  \node[note,text=BookTeal] at (59,68.5) {EMA};

  \node[font=\figurefont\bfseries\footnotesize,text=BookInk,anchor=west] at (5,37) {(c) 对称匹配：VICReg};
  \node[note,anchor=east] at (163,37) {两路共享编码器与扩展头；无停止梯度};
  \node[input] (v-va) at (10,22) {$\bm{v}_A$};
  \node[compute,minimum width=28mm] (v-fa) at (60,22) {$h_{\bm{\phi}}\circ f_{\bm{\theta}}$};
  \node[output,minimum width=14mm] (v-za) at (100,22) {$\bm{Z}$};
  \node[input] (v-vb) at (10,4) {$\bm{v}_B$};
  \node[compute,minimum width=28mm] (v-fb) at (60,4) {$h_{\bm{\phi}}\circ f_{\bm{\theta}}$};
  \node[output,minimum width=14mm] (v-zb) at (100,4) {$\bm{Z}'$};
  \node[loss,minimum width=34mm] (v-loss) at (148,13) {匹配$s(\bm{Z},\bm{Z}')$\\方差$v(\bm{Z}),v(\bm{Z}')$\\协方差$c(\bm{Z}),c(\bm{Z}')$};
  \draw[forward] (v-va.east)--(v-fa.west);
  \draw[forward] (v-fa.east)--(v-za.west);
  \draw[forward] (v-vb.east)--(v-fb.west);
  \draw[forward] (v-fb.east)--(v-zb.west);
  \draw[NNWire,line width=1pt,rounded corners=1mm] (v-za.east) -- (123,22) -- (123,13);
  \draw[NNWire,line width=1pt,rounded corners=1mm] (v-zb.east) -- (123,4) -- (123,13);
  \draw[forward] (123,13)--(v-loss.west);
  \draw[gradient] (v-loss.north) -- (148,29) -- (60,29) -- (v-fa.north);
  \draw[gradient] (100,29)--(v-za.north);
  \draw[gradient] (v-loss.south) -- (148,-3) -- (60,-3) -- (v-fb.south);
  \draw[gradient] (100,-3)--(v-zb.south);

  \draw[forward] (3,-17) -- (16,-17);
  \node[note,anchor=west] at (18,-17) {前向值};
  \draw[gradient] (69,-17) -- (56,-17);
  \node[note,anchor=west] at (72,-17) {当前步梯度};
  \draw[ema] (120,-17) -- (133,-17);
  \node[note,anchor=west] at (136,-17) {跨步参数平均};
\end{tikzpicture}

  \caption{三类关系目标的计算与梯度路径。实线为前向值，橙色虚线为当前梯度，青色点划线为跨步EMA。SimCLR两路共享编码器和投影器，损失还使用批内对照；BYOL只画$A\to B$一项，目标侧在sg处停止梯度，只有编码器与投影器参与EMA，交换视图另算一项；VICReg两路共享$f_{\bm{\theta}}$与扩展头$h_{\bm{\phi}}$，匹配项及各路统计项都向相应分支回传。}
  \label{fig:data-derived-knowledge-relation-objectives}
\end{figure*}
\FloatBarrier

表~\ref{tab:data-derived-knowledge-relation-objectives}汇总已经建立的比较对象与更新规则。SimCLR还需要负对照，其余三种方法只利用正对应；但后者既可以使用独立目标分支，也可以让两路共同求导。“无显式负例”“有教师”“停止梯度”和“对称化”不能互相推出。

\begin{table*}[htbp]
  \centering
  \renewcommand{\arraystretch}{1.15}
  \begin{tabularx}{\linewidth}{@{}>{\RaggedRight\setlength{\parindent}{0pt}\arraybackslash}p{22mm}*{3}{>{\RaggedRight\setlength{\parindent}{0pt}\arraybackslash}X}@{}}
    \toprule
    方法 & 比较对象 & 当前梯度去向 & 跨步变化 \\
    \midrule
    SimCLR & 投影表示的相对相似度 & 两路编码器与投影器 & 共享参数更新后重算表示 \\
    BYOL & 在线预测与目标投影 & 每项仅回传在线分支 & 目标编码器与投影器采用EMA \\
    DINO & 学生与教师的输出分布 & 每项仅回传学生分支 & 教师参数与输出中心分别移动平均 \\
    VICReg & 嵌入匹配与批内统计 & 两路编码器与扩展头 & 共享参数更新后重算表示及统计量 \\
    \bottomrule
  \end{tabularx}
  \caption{关系依据相同时仍可选择不同训练目标。BYOL在线侧另有预测器，DINO没有该预测器；VICReg没有独立教师。各方法的退化控制与成本见正文，不能由某一个组件推断完整方法的保证。}
  \label{tab:data-derived-knowledge-relation-objectives}
\end{table*}
\FloatBarrier

这些方法可以接入同一个最小下游评价流程：冻结各自选定的编码器，读取特征$\bm{h}=f(\bm{x})$，再用独立的任务标签拟合线性分类头。这样比较的是预训练特征支持任务的程度，而不是训练头对自身目标的拟合程度。若同时更新编码器，则属于整体微调，应另报评价协议；不能把投影向量、BYOL预测向量和DINO的$K$维概率都当成同一种下游特征。

\subsection{目标更新与表示退化的控制}
\label{sec:data-derived-knowledge-target-update}

表示目标会随参数变化，因此必须把一个训练步拆成四件事：前向计算目标、按允许路径反向传播、更新接受梯度的参数、准备下一步目标。把这几件事写成一句“教师随学生更新”，会隐藏停止梯度究竟阻断了什么。

另一方面，小匹配误差不等于有信息的表示。退化控制要问完整目标还施加了哪些区分、方差或去相关要求，以及哪些退化状态受到惩罚、哪些仍可能成为驻点。更新机制与退化控制相关，却不能互相替代。

\subsubsection{目标表示的更新}
\label{sec:data-derived-knowledge-target-representation-update}

共享参数的两路表示在同一次优化后一起变化，例如SimCLR和常见VICReg实现。独立目标分支则可以固定、周期性复制或移动平均。固定教师在全部训练期间参数不变；停止梯度只让某条路径在当前求导中视为常量；不计算目标表示则连前向值也不存在。三种说法必须分开。

表~\ref{tab:data-derived-knowledge-gradient-and-update}按时点追踪BYOL式更新。以下数值是教学构造：某个参与EMA的标量在线参数$\theta^{(t)}=2.00$，当前梯度为$0.50$，学习率$\eta=0.20$，目标参数$\xi^{(t)}=1.60$，EMA系数$m=0.90$。在线更新为$2.00-0.20\times0.50=1.90$，随后目标更新为$0.90\times1.60+0.10\times1.90=1.63$；预测器参数不参与这一步平均。

\begin{table*}[htbp]
  \centering
  \renewcommand{\arraystretch}{1.15}
  \begin{tabularx}{\linewidth}{@{}>{\RaggedRight\setlength{\parindent}{0pt}\arraybackslash}p{25mm}*{3}{>{\RaggedRight\setlength{\parindent}{0pt}\arraybackslash}X}@{}}
    \toprule
    时点 & 在线分支 & 目标分支 & 本步操作 \\
    \midrule
    1. 前向 & $\bm{\theta}^{(t)}$产生预测 & $\bm{\xi}^{(t)}$产生目标 & 比较预测与停止梯度的目标值 \\
    2. 反向 & 接收梯度；示例为$0.50$ & 不接收损失梯度 & 只沿在线路径累计梯度 \\
    3. 在线更新 & 优化器更新；示例为$1.90$ & 仍为$\bm{\xi}^{(t)}$ & 先完成在线参数更新 \\
    4. 目标更新 & 保持$\bm{\theta}^{(t+1)}$ & EMA更新；示例为$1.63$ & 用式~\eqref{eq:data-derived-knowledge-byol-ema}平均对应参数 \\
    5. 下一步 & 新参数产生新预测 & 新参数重新计算目标 & 进入$t+1$步前向计算 \\
    \bottomrule
  \end{tabularx}
  \caption{当前步梯度与跨步目标更新。固定教师会省略第4行，目标参数保持不变；共享参数方案没有独立$\bm{\xi}$，两路随同一在线更新改变；VICReg等两路共同反向传播的方法按图~\ref{fig:data-derived-knowledge-relation-objectives}处理。}
  \label{tab:data-derived-knowledge-gradient-and-update}
\end{table*}
\FloatBarrier

EMA系数$m$控制目标混合新旧参数的速度。$m$较小，目标紧跟在线网络，变化快；$m$接近$1$，目标更平滑，却可能明显滞后。具体日程还与优化器步长和训练长度共同作用，不能孤立比较$m$。多一套目标参数也会增加模型存储和前向计算。

目标更新顺序同样重要。式~\eqref{eq:data-derived-knowledge-byol-ema}使用已经完成在线更新的$\bm{\theta}^{(t+1)}$；若实现改用$\bm{\theta}^{(t)}$，就多引入一步滞后。分布式训练还需确认各设备在EMA前是否已经同步在线参数，否则“同一目标网络”可能在设备间并不一致。

DINO除教师参数EMA外，还维护教师输出中心的移动平均。前者改变生成特征与输出的网络，后者改变教师输出居中的基准；二者保存和更新的对象不同。

\subsubsection{表示退化的控制}
\label{sec:data-derived-knowledge-representation-degeneration}

\term{表示坍塌}（representation collapse）最直接的形式是所有样本输出同一向量。若目标只要求正对应接近，这个常数解能让匹配误差为零。还有较弱的\term{维度退化}：样本之间有变化，但多个坐标重复表达同一方向，表示的有效维数远小于名义维数。

图~\ref{fig:data-derived-knowledge-representation-degeneration}用四个教学样本展示三种情况。每个样本的两视图都很接近，所以单看配对误差，三个面板都可能“成功”；跨样本方差和坐标协方差却能区分常数输出与单方向变化。右侧分散只表示这两个示意维度没有前两种简单退化，不表示形成了真实类别。

\begin{figure*}[!t]
  \centering
  \begin{tikzpicture}[
  x=1mm,y=1mm,
  font=\figurefont\scriptsize,
  panel/.style={draw=BookRule,fill=NNPanel,line width=.85pt,rounded corners=1.5mm},
  axis/.style={knowledge arrow,draw=BookMuted,line width=.8pt},
  viewone/.style args={#1}{circle,fill=#1,draw=#1,minimum size=3.2mm,inner sep=0pt},
  viewtwo/.style args={#1}{circle,fill=white,draw=#1,line width=1pt,minimum size=3.8mm,inner sep=0pt},
  samplelabel/.style={font=\figurefont\scriptsize,text=BookInk},
  note/.style={font=\figurefont\scriptsize,text=BookMuted,align=center}
]
  \draw[panel] (.5,.5) rectangle (51.5,66.5);
  \node[font=\figurefont\bfseries\footnotesize,text=BookInk] at (26,62) {(a) 常数表示};
  \draw[axis] (9,32) -- (43,32) node[right] {$z_1$};
  \draw[axis] (26,20) -- (26,49) node[above] {$z_2$};
  \node[viewtwo=BookAmber,minimum size=7mm] at (26,32) {};
  \node[viewtwo=NNHidden,minimum size=5.8mm] at (26,32) {};
  \node[viewtwo=BookTeal,minimum size=4.6mm] at (26,32) {};
  \node[viewone=BookBlue] at (26,32) {};
  \node[samplelabel,anchor=west] at (30,35) {$A=B=C=D$};
  \node[note] at (26,12) {$\bm{Z}=\bm{Z}'=\begin{bmatrix}0&0\\0&0\\0&0\\0&0\end{bmatrix}$\\配对误差为$0$，各维方差为$0$};

  \draw[panel] (58.5,.5) rectangle (110.5,66.5);
  \node[font=\figurefont\bfseries\footnotesize,text=BookInk] at (84.5,62) {(b) 单方向变化};
  \draw[axis] (67,32) -- (101,32) node[right] {$z_1$};
  \draw[axis] (84,20) -- (84,49) node[above] {$z_2$};
  \draw[BookRule,line width=.6pt,dashed] (72,20) -- (98,46);
  \foreach \x/\y/\sampleid/\samplecolor in {73/21/A/BookBlue,80/28/B/BookTeal,88/36/C/NNHidden,95/43/D/BookAmber}{
    \node[viewtwo=\samplecolor] at ({\x+.7},{\y-.4}) {};
    \node[viewone=\samplecolor] at (\x,\y) {};
    \node[samplelabel,anchor=south east] at ({\x-1},{\y+1}) {\sampleid};
  }
  \node[note] at (84.5,10.5) {$\bm{Z}\approx\bm{Z}'=\begin{bmatrix}-1.5&-1.5\\-.5&-.5\\.5&.5\\1.5&1.5\end{bmatrix}$\\样本有方差，但两个坐标重复同一变化};

  \draw[panel] (117.5,.5) rectangle (168.5,66.5);
  \node[font=\figurefont\bfseries\footnotesize,text=BookInk] at (143,62) {(c) 示意维度中分散};
  \draw[axis] (126,32) -- (160,32) node[right] {$z_1$};
  \draw[axis] (143,20) -- (143,49) node[above] {$z_2$};
  \foreach \x/\y/\sampleid/\samplecolor in {132/22/A/BookBlue,132/42/B/BookTeal,154/22/C/NNHidden,154/42/D/BookAmber}{
    \node[viewtwo=\samplecolor] at ({\x+.7},{\y-.4}) {};
    \node[viewone=\samplecolor] at (\x,\y) {};
    \node[samplelabel,anchor=south east] at ({\x-1},{\y+1}) {\sampleid};
  }
  \node[note] at (143,10.5) {$\bm{Z}\approx\bm{Z}'=\begin{bmatrix}-1&-1\\-1&1\\1&-1\\1&1\end{bmatrix}$\\两维有方差且示例协方差为$0$};

  \node[viewone=BookInk] at (5,-5) {};
  \node[note,anchor=west] at (8,-5) {视图A};
  \node[viewtwo=BookInk] at (37,-5) {};
  \node[note,anchor=west] at (40,-5) {对应视图B};
  \node[note,anchor=west,text=FigureLoss] at (76,-5) {教学构造；分散不等于形成有用类别};
\end{tikzpicture}

  \caption{匹配成功与表示退化可以同时发生。三个面板使用相同坐标范围和样本标记；实心点与空心圈表示同一样本的两个视图。下方矩阵均为手工构造，分别对应全部重合、两个坐标重复同一变化、以及在示意维度中分散的情况。}
  \label{fig:data-derived-knowledge-representation-degeneration}
\end{figure*}

SimCLR在存在有效对照且模型能实现更好区分时，偏好比常数输出更低的相对判别损失；第~\ref{sec:data-derived-knowledge-contrastive}节的代入计算也说明了常数驻点仍可存在。错误负例可能迫使语义相近样本分开，容易区分的背景也可能满足目标。对比项控制的是当前关系集合中的相对排序，不是语义类别本身。

BYOL的预测头、停止梯度、移动平均目标与优化动态共同形成非对称训练。原论文的完整方法没有显式负例，并通过消融研究这些组件\scite{ka-grill2020byol}
不能从“停止梯度”单独推出不会坍塌，也不能把EMA称为沿目标损失梯度更新；式~\eqref{eq:data-derived-knowledge-byol-ema}明确显示它采用另一条参数规则。

DINO针对输出分布区分两种退化：所有样本都接近均匀分布，或者所有样本集中到同一输出维度。教师输出锐化抵抗前者，居中抵抗后者，两者在动量教师条件下共同使用\scite{ka-caron2021dino}
这两种分布退化不同于所有连续向量坐标相等，诊断时不能只画一个重合散点代替。

VICReg把控制要求直接写入损失。在$n\geq2$、$\gamma>\sqrt{\epsilon}$的条件下，方差项对图~\ref{fig:data-derived-knowledge-representation-degeneration}(a)产生惩罚，协方差项对(b)产生惩罚。按图中的教学矩阵及样本方差分母$n-1$计算，(b)的协方差矩阵各元素为$5/3$，所以$c(\bm{Z})=25/9$；(c)的协方差矩阵为$(4/3)\bm{I}_2$，所以$c(\bm{Z})=0$。这些量区分了坐标冗余，却不保证优化能从任意初值避开退化。

这些控制只针对特定的退化形式。较高方差可能来自照明，低协方差也可能只是对背景特征作了去相关处理，分散表示仍可能与目标任务无关。因此至少要分开观察匹配误差、每维方差、协方差或有效秩，以及独立下游任务表现。训练损失降低不能替代这些证据。

\section{基于数据分组的知识获取}
\label{sec:data-derived-knowledge-grouping}

关系方法规定样本对之间怎样比较，分组方法则进一步把许多样本组织成群组，并把组编号或分配向量作为目标。\term{分组目标}是由当前表示、相似性和分组配置共同产生的临时答案。它没有预先给定的任务类别含义，组1在下一轮也不必仍对应当前的组1。

聚类算法的结构假设与求解方法已在第~\ref{chap:clustering}章展开。本节只关注分组怎样进入后续训练：用什么表示形成组，输出硬编号还是软分配，预测器读取什么输入，何时刷新分组，以及刷新后哪些状态必须同步改变。

图~\ref{fig:data-derived-knowledge-grouping-targets}把同一批表示送入三种机制。它们可以共用编码器，也可以组合使用：原型分配本身是一种分组，跨视图交换预测又利用了第~\ref{sec:data-derived-knowledge-relations}节的正对应。分类位置取决于当前要解释的主要知识构造机制。

\begin{figure*}[!t]
  \centering
  \begin{tikzpicture}[
  x=1mm,y=1mm,
  font=\figurefont\scriptsize,
  panel/.style={draw=BookRule,fill=NNPanel,line width=.85pt,rounded corners=1.5mm},
  source/.style={draw=FigureInput,fill=FigureInput!12,line width=1.1pt,minimum height=9mm,minimum width=28mm,align=center,inner sep=1.5mm},
  target/.style={draw=FigureOutput,fill=FigureOutput!10,line width=1.1pt,minimum height=8mm,minimum width=23mm,align=center,inner sep=1.5mm},
  model/.style={draw=FigureModel,fill=FigureModel!12,line width=1.1pt,minimum height=8mm,minimum width=23mm,align=center,inner sep=1.5mm},
  flow/.style={knowledge arrow,draw=NNWire,line width=1.1pt},
  refresh/.style={knowledge arrow,draw=BookTeal,line width=.9pt,dashed,rounded corners=1.5mm},
  edge/.style={draw=NNLine,line width=.8pt},
  warningedge/.style={draw=FigureLoss,line width=1.2pt,dashed},
  point/.style={circle,draw=BookBlue,fill=BookBlue!18,minimum size=3.4mm,inner sep=0pt},
  note/.style={font=\figurefont\scriptsize,text=BookMuted,align=center}
]
  \draw[panel] (.5,-34) rectangle (52.5,34.5);
  \draw[panel] (58.5,-34) rectangle (110.5,34.5);
  \draw[panel] (116.5,-34) rectangle (168.5,34.5);
  \node[source] (batch) at (84.5,57) {同一批未标注样本$\bm{X}$};
  \node[model] (repr) at (84.5,44) {共享表示$\bm{Z}=f_{\bm{\theta}}(\bm{X})$};
  \draw[flow] (batch)--(repr);
  \draw[NNWire,line width=1.1pt] (repr.south) -- (84.5,38);
  \draw[NNWire,line width=1.1pt] (5,38) -- (164,38);

  \node[font=\figurefont\bfseries\footnotesize,text=BookInk] at (26.5,30) {(a) 聚类编号};
  \draw[flow] (5,38) |- (26.5,24);
  \foreach \x/\y in {10/16,14/20,18/14,34/17,39/21,43/14}{\node[point] at (\x,\y) {};}
  \draw[BookBlue,line width=.8pt,dashed] (14,17) ellipse (9mm and 7mm);
  \draw[BookTeal,line width=.8pt,dashed] (39,17) ellipse (9mm and 7mm);
  \node[text=BookBlue] at (14,8) {组1};
  \node[text=BookTeal] at (39,8) {组2};
  \node[target] (cluster-target) at (26.5,-5) {临时目标$c_i$};
  \node[model] (cluster-predict) at (26.5,-24) {预测组号并更新};
  \draw[flow] (26.5,7)--(cluster-target);
  \draw[flow] (cluster-target)--(cluster-predict);
  \draw[refresh] (cluster-predict.east) -| (51,23);
  \node[note,text=BookTeal,rotate=90] at (47,-9) {可选刷新};

  \node[font=\figurefont\bfseries\footnotesize,text=BookInk] at (84.5,30) {(b) 原型软分配};
  \node[note,text=BookInk] (assignments) at (84.5,15) {$\begin{array}{c|cc}
    &\text{样本1}&\text{样本2}\\
    \bm{c}_1&0.690&0.310\\
    \bm{c}_2&0.310&0.690
  \end{array}$};
  \draw[flow] (61,38) |- (assignments.west);
  \node[target] (proto-target) at (84.5,-5) {分配向量$\bm{q}_{s,b}$};
  \node[model] (proto-predict) at (84.5,-24) {让视图$t$预测};
  \draw[flow] (assignments.south)--(proto-target.north);
  \draw[flow] (proto-target)--(proto-predict);
  \draw[refresh] (proto-predict.east) -| (109,23);

  \node[font=\figurefont\bfseries\footnotesize,text=BookInk] at (142.5,30) {(c) 邻域／图分组};
  \draw[flow] (119,38) |- (142.5,26.5);
  \node[point,minimum size=4.3mm] (g1) at (125,17) {1};
  \node[point,minimum size=4.3mm] (g2) at (132,23) {2};
  \node[point,minimum size=4.3mm] (g3) at (136,12) {3};
  \node[point,minimum size=4.3mm] (g4) at (150,12) {4};
  \node[point,minimum size=4.3mm] (g5) at (156,23) {5};
  \node[point,minimum size=4.3mm] (g6) at (163,17) {6};
  \draw[edge] (g1)--(g2);
  \draw[edge] (g2)--(g3);
  \draw[edge] (g3)--(g1);
  \draw[edge] (g4)--(g5);
  \draw[edge] (g5)--(g6);
  \draw[edge] (g6)--(g4);
  \draw[warningedge] (g3.east)-- node[above=1mm,text=FigureLoss] {桥接} (g4.west);
  \node[target,text width=38mm,inner sep=1mm] (graph-target) at (142.5,-5) {前：$(1,1,1,2,2,2)$\\后：$(1,1,1,1,1,1)$};
  \node[model] (graph-predict) at (142.5,-24) {预测目标并更新};
  \draw[flow] (142.5,7)--(graph-target);
  \draw[flow] (graph-target)--(graph-predict);
  \draw[refresh] (graph-predict.east) -| (167,23);

  \node[note,anchor=north west,text width=166mm,align=left] at (1,-38) {原型面板只展开两个样本，列内权重和为1；图面板展开六个节点。数值均为教学构造，分组和原型索引不具有预设物种含义。};
\end{tikzpicture}

  \caption{分组结构怎样变成训练目标。聚类产生临时组号；原型面板用两个样本示意分配向量，数值对应第~\ref{sec:data-derived-knowledge-prototype-targets}节的教学算例；图面板的节点1至6对应第~\ref{sec:data-derived-knowledge-graph-targets}节的组号计算，赭色虚线是待加入的错误桥接边$(3,4)$。青色回线表示更新表示后可重算目标，不是通过分组求解器反向传播。}
  \label{fig:data-derived-knowledge-grouping-targets}
\end{figure*}

表~\ref{tab:data-derived-knowledge-grouping}进一步比较三种机制产生的目标、统计范围和刷新要求。分组一旦更新，目标编号、原型或图关系及其下游接口必须同步变化，不能只替换其中一项。

\begin{table*}[htbp]
  \centering
  \renewcommand{\arraystretch}{1.15}
  \begin{tabularx}{\linewidth}{@{}>{\RaggedRight\setlength{\parindent}{0pt}\arraybackslash}p{25mm}*{3}{>{\RaggedRight\setlength{\parindent}{0pt}\arraybackslash}X}@{}}
    \toprule
    分组方式 & 输出目标 & 统计范围 & 刷新时同步对象 \\
    \midrule
    聚类分组 & 硬簇编号$c_i$ & 常需全数据或较大样本 & 簇中心、样本目标、预测头与编号对应 \\
    原型分配 & 硬编号或软分配$\bm{q}_i$ & 小批量，可辅以表示队列 & 原型、分配结果与队列 \\
    邻域／图分组 & 组号或组内关系 & 局部近邻或全局图 & 邻接表、分组结果与训练目标 \\
    \bottomrule
  \end{tabularx}
  \caption{分组目标的生成与维护要求。三种机制并不互斥；编号和分配向量都是算法构造，不是预先给定的类别答案。具体配置及失效条件在各小节展开。}
  \label{tab:data-derived-knowledge-grouping}
\end{table*}

% 首段的DeepCluster边注从段末开始，避免贴近纸张底边。
\Needspace{65mm}
\subsection{基于聚类的分组目标构造}
\label{sec:data-derived-knowledge-cluster-targets}

最直接的做法是先提取每个样本的表示$\bm{z}_i=f_{\bm{\theta}}(\bm{x}_i)$，运行聚类得到$c_i\in\{1,\ldots,K\}$，再训练模型根据输入预测$c_i$。DeepCluster给出了这一交替机制的代表实现：用$k$均值对网络特征分组，再以簇分配监督网络更新\scite{knowledge-acquisition-deepcluster2018}

簇编号把无标签样本转成了一个分类式目标，但它的含义完全由当前表示和聚类配置决定。下面的流程是对这类交替训练的抽象，不把编号对齐、刷新策略和失败边界都归因于某一篇论文。

固定分组只运行一次聚类，之后一直预测同一批$c_i$。周期性分组则交替执行
\[
  \begin{aligned}
    \{\bm{z}_i^{(r)}=f_{\bm{\theta}^{(r)}}(\bm{x}_i)\}_{i=1}^{n}
    &\ \longrightarrow\
    \bm{c}^{(r)}
    =\operatorname{Cluster}(\{\bm{z}_i^{(r)}\}_{i=1}^{n}),\\
    \mathcal{D}^{(r)}=\{(\bm{x}_i,c_i^{(r)})\}_{i=1}^{n}
    &\ \longrightarrow\
    \bm{\theta}^{(r+1)}=\operatorname{Train}(\mathcal{D}^{(r)}),
  \end{aligned}
\]
其中$\bm{c}^{(r)}=(c_i^{(r)})_{i=1}^{n}$，$r$是分组轮次而不是梯度步。前一轮训练改变表示，下一轮聚类便可能改变目标；频繁刷新更及时，也增加全量编码、聚类和目标存储成本。

簇编号没有天然顺序。即使两轮几何分组完全相同，算法也可能把原来的组1改名为组3。若预测头跨轮保留，就要按样本重叠或中心匹配对齐编号，或者重新初始化相应输出；仅按数字相同继承参数没有依据。群组发生分裂和合并时，一一匹配甚至不存在。

未标注花卉照片可能按背景颜色聚成组，也可能按花瓣形态聚成组。两种划分都可以降低各自的聚类目标，却对物种识别提供不同信息。簇数、特征尺度、异常点、组间不平衡和初始表示偏向都会影响结果，错误分组还会被后续预测器强化。

因此，簇编号只能称为临时目标。只有在额外标签、领域解释、独立数据稳定性或下游用途提供证据后，才能把某个簇解释为任务类别。聚类目标下降和预测簇号准确，只说明模型复现了当前分组。

\subsection{基于原型的分配目标构造}
\label{sec:data-derived-knowledge-prototype-targets}

\term{原型}（prototype）是一组可更新的代表向量$\bm{c}_1,\ldots,\bm{c}_K$。给定表示$\bm{z}_i$，可以选择最近原型形成硬编号，也可以保留到多个原型的软权重$\bm{q}_i\in\mathbb{R}^{K}$。软分配表达当前样本对多个代表的相对支持，却不自动具有类别后验概率的语义。

SwAV把原型分配与跨视图关系结合。对同一图像的两个视图$t,s$，归一化表示为$\bm{z}_t,\bm{z}_s$，相应分配为$\bm{q}_t,\bm{q}_s$。用表示$\bm{z}$预测原型$k$的概率为
\[
  p_k(\bm{z})
  =\frac{\exp(\bm{z}^{\top}\bm{c}_k/\tau)}
  {\sum_{r=1}^{K}\exp(\bm{z}^{\top}\bm{c}_r/\tau)}.
\]
交换预测损失让一个视图预测另一个视图的分配：
\begin{equation}
  \mathcal{L}_{\mathrm{swap}}
  =-\sum_{k=1}^{K}q_{s,k}\log p_k(\bm{z}_t)
   -\sum_{k=1}^{K}q_{t,k}\log p_k(\bm{z}_s).
  \label{eq:data-derived-knowledge-swav-swap}
\end{equation}
因此，$\bm{q}_s$来自视图$s$，却监督视图$t$的原型概率，反向亦然\scite{ka-caron2020swav}

分配不是独立地把每个样本送到最近原型。对一个含$B$个表示的批次，写$\bm{Z}=[\bm{z}_1,\ldots,\bm{z}_B]\in\mathbb{R}^{d\times B}$、$\bm{C}=[\bm{c}_1,\ldots,\bm{c}_K]\in\mathbb{R}^{d\times K}$，SwAV求解带熵正则的分配：
\begin{align}
  \max_{\bm{Q}\in\mathcal{Q}}\quad&
  \operatorname{Tr}(\bm{Q}^{\top}\bm{C}^{\top}\bm{Z})
  +\varepsilon H(\bm{Q}),\\
  \mathcal{Q}
  ={}&
  \left\{
    \bm{Q}\in\mathbb{R}_{+}^{K\times B}
    \ \middle|\
    \begin{aligned}
      \bm{Q}\bm{1}_B&=\frac{1}{K}\bm{1}_K,\\
      \bm{Q}^{\top}\bm{1}_K&=\frac{1}{B}\bm{1}_B
    \end{aligned}
  \right\}.
  \label{eq:data-derived-knowledge-swav-assignment}
\end{align}
这里$H(\bm{Q})=-\sum_{k=1}^{K}\sum_{b=1}^{B}Q_{k,b}\log Q_{k,b}$，$\varepsilon>0$控制熵正则强度，$\bm{1}_B$和$\bm{1}_K$分别是$B$维和$K$维全一向量。约束要求每行之和为$1/K$、每列之和为$1/B$，总权重为$1$。它们在数学上是严格等式，数值迭代只会近似满足。由于$\bm{1}_K^{\top}\bm{Q}_{:,b}=1/B$，第$b$列必须在进入交叉熵前归一化为
\begin{equation}
  \bm{q}_b
  =
  \frac{\bm{Q}_{:,b}}
       {\bm{1}_K^{\top}\bm{Q}_{:,b}}
  =
  B\bm{Q}_{:,b},
  \qquad
  \bm{1}_K^{\top}\bm{q}_b=1.
  \label{eq:data-derived-knowledge-swav-code-normalization}
\end{equation}
式~\eqref{eq:data-derived-knowledge-swav-swap}中的$\bm{q}_t,\bm{q}_s$固定表示这种归一化的\term{分配向量}，原论文称为code；这里不是程序代码，也不是和为$1/B$的原始列。这样，交换预测损失中的目标才是原型上的概率分布，不会额外缩小$1/B$。

如何从相似度得到$\bm{Q}$？令分数矩阵$\bm{F}=\bm{C}^{\top}\bm{Z}$，先逐元素指数化并归一化总和：
\[
  Q_{k,b}^{(0)}
  =\frac{\exp(F_{k,b}/\varepsilon)}
  {\sum_{r=1}^{K}\sum_{a=1}^{B}\exp(F_{r,a}/\varepsilon)}.
\]
Sinkhorn--Knopp迭代交替调整行和与列和。在第$t$轮中，先得到行缩放结果$\bm{R}^{(t)}$，再得到列缩放结果$\bm{Q}^{(t)}$：
\begin{equation}
  \begin{aligned}
    R_{k,b}^{(t)}
    &=\frac{Q_{k,b}^{(t-1)}}{K\sum_{a=1}^{B}Q_{k,a}^{(t-1)}},\\
    Q_{k,b}^{(t)}
    &=\frac{R_{k,b}^{(t)}}{B\sum_{r=1}^{K}R_{r,b}^{(t)}}.
  \end{aligned}
  \label{eq:data-derived-knowledge-sinkhorn-scaling}
\end{equation}
一次列缩放可能再次改变行和，所以不能只做一次就默认所有边际已经达标。这里以各行偏离$1/K$及各列偏离$1/B$的最大绝对值小于$10^{-10}$为停止条件，并总在列缩放后检查；此时列和在数值精度内为$1/B$，可以使用$\bm{q}_b=B\bm{Q}_{:,b}$。若工程实现固定只做少量轮次，应另报尚存的行边际误差。

下面用完全人为构造的两样本、两原型例子复算全过程，不代表训练实验。取$K=B=2$，单位原型$\bm{c}_1=(1,0)^{\top}$、$\bm{c}_2=(0,1)^{\top}$，视图$s$的单位表示为$\bm{z}_{s,1}=(1,0)^{\top}$、$\bm{z}_{s,2}=(0.8,0.6)^{\top}$，并取$\varepsilon=\tau=0.5$。分数和初始化矩阵为
\[
  \bm{F}_s=\begin{bmatrix}1&0.8\\0&0.6\end{bmatrix},
  \qquad
  \bm{Q}^{(0)}\approx
  \begin{bmatrix}0.443462&0.297262\\0.060016&0.199260\end{bmatrix}.
\]
首轮行缩放与列缩放分别给出
\[
  \begin{aligned}
    \bm{R}^{(1)}&\approx
    \begin{bmatrix}0.299344&0.200656\\0.115738&0.384262\end{bmatrix},\\
    \bm{Q}^{(1)}&\approx
    \begin{bmatrix}0.360584&0.171525\\0.139416&0.328475\end{bmatrix}.
  \end{aligned}
\]
$\bm{R}^{(1)}$的两列之和约为$0.415081$和$0.584919$，还需分别缩放到$0.5$；$\bm{Q}^{(1)}$的两行之和则约为$0.532109$和$0.467891$。用未舍入值继续迭代，12轮后的最大边际误差约为$1.82\times10^{-11}$，得到
\[
  \bm{Q}_s\approx
  \begin{bmatrix}0.344987&0.155013\\0.155013&0.344987\end{bmatrix},
  \qquad
  \bm{q}_{s,1}\approx\begin{bmatrix}0.689974\\0.310026\end{bmatrix},
  \quad
  \bm{q}_{s,2}\approx\begin{bmatrix}0.310026\\0.689974\end{bmatrix}.
\]
若独立选择最近原型，两个样本都会选择原型1；均衡软分配却让第二个样本更多支持原型2。这说明目标不仅取决于单样本相似度，还取决于整个分配批次。

为接上交换损失，再构造视图$t$的表示$\bm{z}_{t,1}=(0.8,0.6)^{\top}$、$\bm{z}_{t,2}=(1,0)^{\top}$；它们仍按样本编号1、2与视图$s$配对，不能因数值交换而重配样本。对视图$t$独立运行同一分配过程，得到$\bm{q}_{t,1}=\bm{q}_{s,2}$、$\bm{q}_{t,2}=\bm{q}_{s,1}$。样本1的预测概率为
\[
  \bm{p}(\bm{z}_{t,1})\approx(0.598688,0.401312)^{\top},
  \qquad
  \bm{p}(\bm{z}_{s,1})\approx(0.880797,0.119203)^{\top}.
\]
以自然对数计算，两个方向的交叉熵及其和为
\[
  \begin{aligned}
    -\bm{q}_{s,1}^{\top}\log\bm{p}(\bm{z}_{t,1})&\approx0.637025,\\
    -\bm{q}_{t,1}^{\top}\log\bm{p}(\bm{z}_{s,1})&\approx1.506877,\\
    \mathcal{L}_{\mathrm{swap},1}&\approx2.143902.
  \end{aligned}
\]
本构造中样本2的两项恰好交换，因此按样本平均的双向损失仍为$2.143902$；若还对两个方向取平均，则为$1.071951$。必须先约定求和或平均范围，才能比较损失数值。

当前损失中，Sinkhorn得到的$\bm{q}_t,\bm{q}_s$停止求导；编码器和原型则通过式~\eqref{eq:data-derived-knowledge-swav-swap}中$p_k(\bm{z})$一侧的梯度更新。优化器完成更新后，原方法把每个原型列投影回单位球，即令$\bm{c}_k\leftarrow\bm{c}_k/\lVert\bm{c}_k\rVert_2$，下一批再用归一化原型计算分数。这个投影是梯度更新后的约束步骤，不是分配目标的反向传播。SwAV没有BYOL式目标网络。若实现让梯度穿过分配求解器，或省略原型归一化，计算路径就与原方法不同。

软均衡约束始终可行：对任意正整数$K,B$，取$Q_{k,b}=1/(KB)$就严格满足全部行列边际，即使所有表示完全相同也如此。当$K>1$时，均衡排除了全部分配权重集中到同一原型，却没有排除每个样本都均匀分配到所有原型。困难在于得到有信息的近硬分配：有效多样性不足时，均衡可能迫使相似样本分摊到多个原型，使分配更软，或与相似度支持的分组冲突。原论文用前批表示队列扩大小批量的分配统计范围，补充可能缺少的多样性，不是修复一个空的软约束集；队列也带来表示过时问题。熵系数$\varepsilon$过大还会进一步偏好均匀分配。

原型数量、相似度温度、分配熵、均衡强度和小批量组成共同决定目标。比较最近原型硬编号与软分配时，应固定原型和表示；硬化会删除不确定性，软分配则增加存储与交叉熵计算。原型编号与聚类编号一样，没有预设物种含义。

SwAV同时使用正对应和分组目标：同一图像的视图关系由数据变换建立，Sinkhorn分配向量来自本节的分组机制。把它只称为“非对比匹配”会漏掉原型分配约束，把它只称为“聚类”又会漏掉交换视图的训练关系。

\subsection{基于邻域与图结构的分组目标构造}
\label{sec:data-derived-knowledge-graph-targets}

中心或原型用少量代表描述群组，邻域方法则直接保留样本间的局部联系。给定表示，可以按$k$近邻、半径阈值或已有关系建立图$\mathcal{G}=(V,E)$，其中顶点对应样本，边表示当前规则认定的邻近。目标可以是局部邻居集合、连通分量编号，也可以来自谱划分或社区发现。

局部邻近不自动具有传递性。样本$A$接近$B$、$B$接近$C$，并不保证$A$与$C$应共享语义；一个桥接点便可能把两个密集群组连成同一分量。图~\ref{fig:data-derived-knowledge-grouping-targets}(c)用赭色边标出了这种风险。若分组依赖连通性，单条错误边的影响可以远大于其局部权重。

固定一种规则就能看见目标怎样改变。取图中的六个教学节点$V=\{1,\ldots,6\}$，以对称相似度不低于$0.8$作为无向连边条件，再用连通分量生成组号，按各分量最小节点编号的顺序命名。设两个三元组$\{1,2,3\}$、$\{4,5,6\}$内部两两相似度均为$0.9$，跨组相似度均为$0.2$，于是
\[
  \begin{aligned}
    E_{\mathrm{before}}
    &=\{(1,2),(1,3),(2,3),(4,5),(4,6),(5,6)\},\\
    \bm{c}_{\mathrm{before}}&=(1,1,1,2,2,2).
  \end{aligned}
\]
若一次错误匹配把节点3与4的相似度改成$0.85$，其他分数不变，新增边$(3,4)$就使所有节点连通，得到$\bm{c}_{\mathrm{after}}=(1,1,1,1,1,1)$。若将同组的不同节点两两转为正对应，桥接前只有两个组内的6对，桥接后变成全部15对，增加了9条跨原组关系。例如训练记录$(\bm{x}_4,c_4)$的目标由2变为1，节点1与6也被加入正对应，尽管它们的直接相似度仍只有$0.2$。这些分数和组号都是教学构造；它们展示的是连通分组的传递作用，不是六个对象的真实语义。

构造时应明确节点、度量、邻居数或阈值、边是否有向、边权含义和最终分组规则。$k$固定时，稀疏区域可能连接很远的样本；半径固定时，密集区域又可能形成过多边。互近邻、局部尺度或密度规则能改变这些行为，但没有脱离数据分布的统一最优选择。

图分组刷新需要重新编码、近邻检索、建边和划分。用直接全对比较构建精确近邻时，需要$O(n^2)$次距离计算；只有显式保存完整两两距离矩阵时，存储才是$O(n^2)$，只保留每个样本的$k$个近邻时，最终邻接表为$O(nk)$，检索索引另计。近似检索降低计算成本，却会漏边或加错边。重新构图后，节点组号同样可能置换、分裂或合并，应同步刷新训练目标并保留未确定关系。

若群组随后被转换为正对应，可按第~\ref{sec:data-derived-knowledge-relations}节使用；此时正对应的可信程度继承分组假设。它与利用已有类别标签传播答案不同：本节先从数据关联发现群组，没有预先拥有类别语义。只有额外证据才能把图分组解释成任务类别。

邻域与图方法特别容易形成反馈回路：当前表示决定近邻，近邻目标又让这些样本更接近。稳定的错误连接会被不断强化。保留边的来源、距离裕量和跨轮稳定性，并在独立用途上验证，比只观察图损失下降更重要。

\section{本章小结}
\label{sec:data-derived-knowledge-summary}

没有现成任务标签时，数据仍能提供三类训练依据。保留词元、像素或未来片段，得到的是在模型之外已经确定的观测目标；建立同一对象、自然配对和负对照关系，得到的是表示之间应怎样比较的约束；从表示中发现簇、原型或图群组，得到的是由分组算法产生的临时目标。三类依据可以联合使用，但来源不能因目标函数都写成预测而被抹平。

观测目标的核心检查是输入是否真的看不到答案，以及损失位置是否与构造一致。关系目标还要追踪比较对象和梯度路径：SimCLR用正对应与候选对照作相对判别，BYOL用在线分支预测停止梯度的EMA目标，DINO匹配居中、锐化后的教师分布，VICReg让匹配与两路统计项共同反向传播。停止梯度只描述当前步，目标分支在下一步可以因共享参数、复制或移动平均而改变。

分组目标则需要区分组号与任务类别。聚类编号可能跨轮置换，原型软分配依赖温度和均衡约束，图分组会放大错误桥接。SwAV说明这些维度可以交叉：分配来自原型分组，交换预测又使用跨视图正对应。组合方法仍须说明每个目标的来源、更新顺序和兼容条件。

减少人工任务标签不意味着没有人为选择，也不意味着构造出的目标必然有用。遮蔽比例、视图变换、负对照采样、教师更新、统计约束、原型数量和建图尺度都表达了假设并消耗计算。评价时应分别检查构造条件是否成立、表示是否保留足够差异、这些差异是否服务下游任务；第~\ref{chap:knowledge-fusion}章将讨论怎样组合多种获取途径，第~\ref{chap:knowledge-quality-assessment-correction}章再把这些证据组织成可复核的协议。
